\section{Unidad 1}

\subsection{¿Cuántos electrones máximo puede existir en la cuarta capa de un átomo?}
\[
Ne = 2n^2 = 2(4)^4 = 2(16) = 32
\]

\subsection{¿Un material aislante nunca podrá conducir corriente eléctrica?}

\subsection{¿Qué relación encuentra entre las capas de valencia del silicio y el germanio, respecto a los electrones de su última capa?}

Los electrones de valencia del germanio se encuentran a niveles de energía más altos que aquellos en el silicio y, por consiguiente, requieren una cantidad de energía adicional más pequeña para escaparse del átomo.

\subsection{¿En qué condición cree que se encuentra este átomo de silicio para tener cero excitación?}

Esta condición de excitación nula se supone ausencia de energía externa absoluta como puede ser el mismo calor. Es decir, el átomo se encuentra a la temperatura de 0°Kelvin.

\subsection{¿La circulación de corriente en el cobre tiene el mismo comportamiento que el silicio?}

Los átomos de cobre forman un tipo de cristal diferente en el que los átomos no están enlazados covalentemente entre sí, sino que se componen de un “mar” de núcleos de iones positivos, los cuales son átomos sin sus electrones de valencia. Los electrones de valencia están enlazados a los iones positivos, lo que mantiene a los iones positivos juntos y les permite formar el enlace metálico. Los electrones de valencia no pertenecen a un átomo dado, sino al cristal en conjunto. Debido a que los electrones de valencia en el cobre se mueven libremente, la aplicación de un voltaje produce corriente. Existe sólo un tipo de corriente –el movimiento de electrones libres– porque no existen “huecos” en la estructura cristalina metálica.

\subsection{¿Cómo reaccionarán estos portadores en el momento de la unión de ambos tipos de semiconductores?}

\subsection{¿Por qué se denomina región de empobrecimiento?}

El término empobrecimiento se refiere al hecho de que la región cercana a la unión pn se queda sin portadores de carga (electrones y huecos) debido a la difusión a través de la unión.

\subsection{¿Por qué se establece el equilibrio en la región de empobrecimiento?}

\subsection{¿De qué depende el potencial de barrera en un diodo?}

El potencial de barrera de una unión pn depende de varios factores, incluido el tipo de material semiconductor, la cantidad de dopado y la temperatura. El potencial de barrera típico es aproximadamente de 0.7V para el silicio y de 0.3V para el germanio a 25°C.

\subsection{¿Cuáles son las condiciones que debe cumplir el circuito para que el diodo conduzca corriente?}

Es necesario que lado negativo de VPOLARIZACIÓN está conectado a la región n del diodo y el lado positivo está conectado a la región p. Además, que el voltaje de polarización, VPOLARIZACIÓN, debe ser más grande que el potencial de barrera.

\subsection{¿Cómo es que en polarización directa hay flujo de cargas (corriente)?}

Como las cargas diferentes se atraen, el lado positivo de la fuente de voltaje de polarización atrae los electrones de valencia hacia el extremo izquierdo de la región p. Los huecos en la región p proporcionan el medio o “ruta” para que estos electrones de valencia se desplacen hacia la región p. Los electrones de valencia se desplazan de un hueco al siguiente hacia la izquierda. Los huecos, que son portadores mayoritarios en la región p, efectivamente (no en realidad) se desplazan a la derecha hacia la unión.

\subsection{¿Esta corriente inversa (insignificante) se mantendrá para cualquier polarización inversa?}

\subsection{¿Cómo limitamos la corriente inversa para evitar al ruptura?}

La mayoría de los diodos no son operados en condición de ruptura en inversa, pero si se limita la corriente (por ejemplo mediante la adición de un resistor limitador en serie), el diodo no sufre daños permanentes.

\subsection{¿Qué sucederá con valores negativos de VD?}
\[
ID = ISe^{\frac{VD}{nVT}} - IS
\]

Con valores negativos de VD el término exponencial se reduce con rapidez por debajo del nivel de I y la ecuación resultante para ID es

\[
ID \approx -IS
\]

